\section{Harmonic deformation coordinates}\label{sec:coordinates}

Give $L_i$ its $\omega_0$-Hermitian--Einstein metric and $E_0$ the
direct-sum metric $h_0$. The entire parallel trace-free algebra is the
determinant-normalized copy of $\mathfrak{gl}_2$. All complements and
Sobolev spaces below are fixed at this reference metric; the parameters
vary on one sufficiently small common neighborhood.

\subsection{The harmonic pairing}

\begin{lemma}\label{lem:pairing}
The two cross-$\Ext^1$ spaces have orthonormal Serre-dual bases for
which the harmonic norm on $\Ext^1(E_0,E_0)$ is the Euclidean norm
on $\cV$ and the quadratic obstruction has zero locus $BA=0$.
\end{lemma}

\begin{proof}
Put $L=L_1^*\otimes L_2$, with its induced metric, and give $L^*$
the dual metric. Normalize the parallel holomorphic two-form $\Omega$
so that in a unitary coframe
\[
\omega_0=i(\theta^1\wedge\bar\theta^1+\theta^2\wedge\bar\theta^2),
\qquad \Omega=\theta^1\wedge\theta^2,\qquad
\Omega\wedge\bar\Omega=dV_0.
\]
For a local unitary frame $e$ of $L$, define
\[
J_L\bigl((a_1\bar\theta^1+a_2\bar\theta^2)e\bigr)
=(-\bar a_2\bar\theta^1+\bar a_1\bar\theta^2)e^*.
\]
This is frame-independent and antiunitary, with $J_{L^*}J_L=-I$.
For the bilinear Serre pairing
$Q_L(\alpha,\gamma)=\int_X\Omega\wedge\alpha\wedge\gamma$,
the pointwise calculation gives
\begin{equation}\label{eq:Serre-unitary}
Q_L(\alpha,J_L\beta)=\langle\alpha,\beta\rangle_{L^2},
\end{equation}
where the Hermitian product is linear in the first argument.

If $\beta$ is harmonic, then for every section $s$ of $L$,
$Q_L(\ddbar_Ls,J_L\beta)=0$. Stokes implies
$\ddbar_{L^*}J_L\beta=0$. For a section $u$ of $L^*$,
\[
\langle\ddbar_{L^*}u,J_L\beta\rangle
=-Q_{L^*}(\ddbar_{L^*}u,\beta)=0
\]
by the same identity and Stokes. Thus $J_L\beta$ is coclosed as well.
An orthonormal basis $\alpha_i$ of $H^{0,1}(L)$ and
$\beta_i=J_L\alpha_i$ therefore have the required duality.

Write $A$ with rows $u_i\in\Hom(\CC^2,\CC)$ and $B$ with
columns $v_i\in\Hom(\CC,\CC^2)$. The cross-term form
\[
x=\begin{pmatrix}
0&\sum_i\beta_i\otimes v_i\\
\sum_i\alpha_i\otimes u_i&0
\end{pmatrix}
\]
has norm $\|A\|^2+\|B\|^2$. If $\eta$ generates
$H^{0,2}(\cO_X)$ with $\int_X\Omega\wedge\eta=1$, the class of
$x\wedge x$ has blocks $-BA\,\eta$ and $\tr(BA)\eta$.
Using $\eta=\bar\Omega$ instead multiplies both blocks by
$1/\Vol(X,\omega_0)$. The zero equation is unchanged.
\end{proof}

The parallel normalization is where the Ricci-flat central metric is
used. A general K\"ahler metric need not make the holomorphic two-form
have constant norm. The compatibility with the hyperk\"ahler harmonic
pairing is also described in \cite[arXiv version, Proposition~9.1]{Ver96}.

\subsection{Convergent equivariant quadratic coordinates}

Let $d=\ddbar_0$ on trace-free endomorphism-valued forms, and let
$\mathcal H^1$ be its harmonic degree-one space. The standard harmonic
Kuranishi embedding $\Phi$ satisfies
\begin{equation}\label{eq:harmonic-fixedpoint}
\Phi(b)=b-d^*G_2(\Phi(b)\wedge\Phi(b)),\qquad D\Phi(0)=\iota,
\end{equation}
where $G_2$ is the Green operator in degree two and $\iota$ the harmonic
inclusion. It is holomorphic and equivariant for parallel automorphisms.
Its integrability equation is the harmonic projection of
$\Phi(b)\wedge\Phi(b)$; see \cite[Theorem~1]{BS22} and
\cite[Proposition~4.3]{Del25}. We use reduced analytic spaces.

\begin{proposition}\label{prop:analytic-coordinates}
There is a $G$-equivariant analytic isomorphism from the reduced germ
$\cQ=\{BA=0\}$ to the full reduced harmonic integrability germ.
It has an equivariant ambient holomorphic extension $F$ satisfying
\begin{equation}\label{eq:F-expansion}
F(x)=x+O(|x|^3),\qquad DF(x)=I+O(|x|^2).
\end{equation}
\end{proposition}

\begin{proof}
Formality follows from \cite[Theorem~1.2]{BZ19}, because $E_0$ is
$H$-Gieseker-polystable. The equivariant analytic result
\cite[Proposition~4.4]{AS18}, with its quiver identification, gives
an isomorphism $\psi:Z\to\cQ$ from a Quot slice, tangent to the
identity. We must compare this slice with the \emph{harmonic} family;
formality alone does not supply that comparison.

Write $S$ for the reduced harmonic base. The family on $X\times S$
is holomorphic and analytically semiuniversal in the reduced category
\cite[arXiv version, \S5 and Remark~4]{BS21}. Retain the full automorphism group
$\widetilde G=GL_2\times GL_1$ while acting on this family. For $m$
large, the twist of the central bundle by $mH$ is generated by sections
and has no higher cohomology. After shrinking $S$, proper base change
makes the pushforward of the family a holomorphic vector bundle with
fibre $W=H^0(X,E_0(mH))$, and evaluation remains surjective.
Choose a frame $T_s$ equal to the identity at zero. Average over a
maximal compact $\widetilde K$:
\[
T_s^{\widetilde K}
=\int_{\widetilde K}\rho(k)^{-1}T_{ks}\,k_{\mathcal W}\,dk .
\]
This is a holomorphic equivariant frame after shrinking. The evaluation
quotient defines an analytic equivariant map $q:S\to R$ to the
open algebraic Quot locus, by the analytic Quot theorem
\cite[Th\'eor\`eme~1.1]{Gro61}. Scalars act ineffectively on the base.
Holomorphic compact equivariance complexifies as an identity of germs.

Let $\Gamma=GL(W)$. Choose an $\widetilde G$-invariant complement
$\mathfrak m$ to its stabilizer Lie algebra in $\operatorname{Lie}\Gamma$.
The Luna map $\Gamma\times^{\widetilde G}Z\to R$ is \'etale
\cite[\S III.1]{Lun73}. The exponential local section of
$\Gamma/\widetilde G$ gives an analytic isomorphism of germs
$(u,z)\mapsto e^u z$ from $\mathfrak m\times Z$ to $R$.
This remains true for singular reduced germs. Projecting its inverse
after $q$ yields an equivariant analytic map $h:S\to Z$.
The factor $e^{u(s)}$ identifies the two marked families, so the
Kodaira--Spencer maps give $Dh(0)=I$.

Both reduced tangent spaces are all of $\cV$. The pure-arrow subspaces
$(A,0)$ and $(0,B)$ lie in $\cQ$ and span $\cV$. They also lie in
$S$: their harmonic forms have square zero, and the unique solution
of~\eqref{eq:harmonic-fixedpoint} is the form itself. Reduced
semiuniversality now gives a reverse analytic classifying map
$g:Z\to S$ with $Dg(0)=I$, without requiring $g$ to be equivariant.

An analytic germ self-map with invertible tangent map is an automorphism.
Indeed, in a minimal ambient embedding its lift $H$ is a local
biholomorphism. The induced endomorphism of the local ring is
surjective, since an ambient representative $f$ has preimage
$f\circ H^{-1}$. A surjective endomorphism $\varphi$ of a Noetherian
ring is injective: the chain $\ker\varphi^n$ stabilizes, and lifting
an element of $\ker\varphi$ by a sufficiently high surjective power
then shows it is zero. Applying this to $gh$ and $hg$ proves that
$h$ is an isomorphism. Thus $F=h^{-1}\psi^{-1}$ is an equivariant
analytic comparison. No nonreduced formal miniversal hull, or
convergence of an arbitrary formal series, is used.

Extend the coordinate functions of $F$ to an ambient map $F_0$.
Its derivative is $I$ because the reduced tangent space is all of
$\cV$. Averaging $\int_K k^{-1}F_0(kx)\,dk$ preserves its restriction
to $\cQ$ and makes it equivariant. The ambient inverse-function
theorem applies. The compact automorphism $\diag(-I_2,1)$ acts as
$-I$ on $\cV$, so the averaged map is odd. This proves
\eqref{eq:F-expansion}.
\end{proof}

Set
\begin{equation}\label{eq:reconstruction}
a(x)=\Phi(F(x)),\qquad D_x=d+[a(x),\,\cdot\,],\qquad J_x=Da(x).
\end{equation}
These maps are defined on a full ambient ball. If $C_x\xi$ is the
infinitesimal coefficient action of a parallel complex endomorphism,
equivariance gives the exact identity
\begin{equation}\label{eq:equivariance}
D_x\xi=-J_xC_x\xi.
\end{equation}
For a tangent vector to a smooth piece of $\cQ$, $D_xJ_xv=0$.

\subsection{The corrected real equation}

Let $dV_0=\omega_0^2/2$ and use its fixed pairing on zero-forms.
For the $h_0$-unitary connection associated to a Dolbeault operator
$\partial$, define the Hermitian-valued moment
\begin{equation}\label{eq:infinite-moment}
M_t(\partial)=-\frac{i(F_\partial^0\wedge\omega_t)}{dV_0}.
\end{equation}
The superscript denotes trace-free curvature. Its zero equation is
projective HYM; imposing the determinant line's HE metric separately
gives the fixed-determinant HYM equation. With the affine symplectic
form associated to~\eqref{eq:dolbeault-norm}, integration by parts gives
\begin{equation}\label{eq:infinite-moment-identity}
d\langle M_t,h\rangle(v)
=2\Omega_{\infty,t}(\xi_h,v),\qquad
\xi_h=-D_\partial(ih).
\end{equation}
For example, the quadratic $(1,1)$ curvature is
$-a\wedge a^\dagger-a^\dagger\wedge a$; the $-iF$ moment has
multiplicity-two block $A^*A-BB^*$.

\begin{proposition}\label{prop:corrected-slice}
There is a unique small smooth Hermitian correction $\sigma(t,x)$,
orthogonal to the full parallel trace-free Hermitian algebra, such that
\[
\Psi_t(x)=e^{\sigma(t,x)}\cdot(\ddbar_0+a(x)),\qquad
M_t(\Psi_t(x))
\]
has only its parallel component. The correction is $K$-equivariant and,
in every fixed finite collection of spatial norms,
\begin{equation}\label{eq:sigma-estimate}
\|\sigma(t,x)\|_{C^k}\le C_k(|t|+|x|^2).
\end{equation}
The pulled-back form $\Omega_t=\Psi_t^*\Omega_{\infty,t}$ and centered
finite moment $\rho_t$ satisfy
\begin{equation}\label{eq:form-error}
\|\Omega_{t,x}-\omega^{\rm amb}\|\le C(|t|+|x|^2),\qquad
d\tr(\rho_t h)(v)=2\Omega_t(\xi_h,v),\qquad \rho_t(0)=0.
\end{equation}
In these coordinates the actual HYM equation is precisely
\begin{equation}\label{eq:exact-real-equation}
\rho_t(x)=\tau(t)I_2,\qquad \tau(t)=-28\pi t/3.
\end{equation}
\end{proposition}

\begin{proof}
Project~\eqref{eq:infinite-moment} onto the fixed orthogonal complement
of the parallel Hermitian algebra, after applying a positive gauge
$e^s$. At $(t,x,s)=(0,0,0)$ its derivative in $s$ is a nonzero fixed
multiple of the covariant Laplacian on that complement. The entire
kernel has been removed, so the derivative is an elliptic isomorphism
between fixed high-regularity Sobolev spaces. The implicit-function
theorem and elliptic regularity give the smooth correction and
uniqueness. Its $x$ derivative at zero vanishes by the K\"ahler
identities and harmonicity of $J_0v$. This proves
\eqref{eq:sigma-estimate}; uniqueness gives compact equivariance.
This is also the HYM specialization of the construction in
\cite[Proposition~4.5]{Del25}.

The pullback of the affine symplectic form is closed and has the
finite moment identity. At $(0,0)$ its value is
$\omega^{\rm amb}$, since $D\Psi_0(0)$ is harmonic inclusion.
Invariance under the automorphism acting as $-I$ makes its
coefficient functions even in $x$. Smoothness gives the first
estimate in~\eqref{eq:form-error}. It is symplectic and tames $I$
on a smaller ball. It need not be compatible with $I$.

For $h\in\operatorname{Herm}_2$, the determinant-one lift is
\begin{equation}\label{eq:det-lift}
\widehat h=\diag(S,c),\qquad
S=h-\tfrac13\tr(h)I_2,\quad c=-\tfrac13\tr(h).
\end{equation}
Define $m_t$ by
$\tr(m_t(x)h)=\langle M_t(\Psi_t(x)),\widehat h\rangle_{L^2(dV_0)}$.
At $x=0$ the correction is block scalar by compact invariance, so the
integrated pairing with each parallel block is topological. With
$c_1(L)=i[F_L]/(2\pi)$ and~\eqref{eq:slopes},
\begin{align*}
\tr(m_t(0)h)
&=-2\pi\bigl(\deg_tL_1\,\tr S+\deg_tL_2\,c\bigr)\\
&=-\tfrac{2\pi}{3}\bigl((3-7t)-(3+7t)\bigr)\tr h
=\tfrac{28\pi}{3}t\,\tr h.
\end{align*}
The fixed density cancels in this integral; no volume factor is
present. Set $\rho_t=m_t-m_t(0)$. Its moment identity has the
normalization in~\eqref{eq:form-error}, and vanishing of the actual
moment is exactly~\eqref{eq:exact-real-equation}.
\end{proof}

For $t<0$, the positive real level is the $A$-injective chamber.
This sign also agrees with the slope test for an $A$-only extension:
its subbundle $L_2$ has lower slope than $E$ precisely when $t<0$.
